\begin{table}[htbp]
\centering
\caption{Fiscal-year ends and filing dates in the estimation sample (30-day window)}
\label{tab:timing}
\small
\begin{tabular}{lrr}
\toprule
Month & Fiscal year ends (\%) & 10-K filed (\%) \\
\midrule
Jan & 5.0 & 2.8 \\
Feb & 0.0 & 66.5 \\
Mar & 2.4 & 8.7 \\
Apr & 1.2 & 0.6 \\
May & 1.7 & 2.3 \\
Jun & 5.4 & 1.6 \\
Jul & 1.6 & 1.5 \\
Aug & 1.2 & 5.2 \\
Sep & 5.2 & 1.6 \\
Oct & 2.7 & 1.4 \\
Nov & 1.0 & 5.0 \\
Dec & 72.6 & 2.7 \\
\bottomrule
\end{tabular}
\par\smallskip
\begin{minipage}{\linewidth}\footnotesize
\emph{Notes.} Shares of the 5{,}782 firm-years. Fiscal-year end month: month of the period end, periods ending on 1--7 January counted to December. The median lag between period end and filing is 53 days (10th percentile 38, 90th percentile 59). 145 of the 512 firms have a fiscal year that does not end in December. 73.6\% of the 10-Ks are filed in the calendar year after the fiscal year, and the filings of one fiscal year span a median of 390 calendar days; the filings fall in 169 distinct calendar months.
\end{minipage}
\end{table}
